\section{Fixed-projection edits and finite-event tests}\label{sec:companion}

This appendix gives three results about the choice of instrument and
tests. With fixed projected kernels, a conformal normal form factors
edits by conditional-flux circuits and prime-valuation moves. The
held tests determine precisely which rate ratios and cycle
affinities~\citep{schnakenberg} are invariant. Allowing finite-event comparisons gives a separate
obstruction to fine reachability, even with a single native state.

\subsection{A conformal normal form at fixed projection}

Fix a finite native state space \(X\), a positive rational stationary
distribution \(\mu\), a finite set \(M\) of marks with an involution
\(m\mapsto\bar m\), and a surjection \(b:X\to B\) onto a finite set.
For each mark fix a rational stochastic matrix \(Q_m\) on \(B\).
A native law has positive rational rates \(\theta_m\) summing to one
and rational conditional stochastic kernels \(K_m\) on \(X\). Its
marked transition probabilities are \(P_m(x,y)=\theta_m K_m(x,y)\). Set
\(C_m(x,y)=\mu(x)K_m(x,y)\) and require
\begin{equation}\label{eq:fixed-flux}
 C_{\bar m}=C_m^{\mathsf T},\qquad
 \sum_{y:\,b(y)=v}C_m(x,y)=\mu(x)Q_m(b(x),v).
\end{equation}
The transpose condition makes both margins of every \(C_m\) equal to
\(\mu\). A positive marked edge has stationary flux ratio
\begin{equation}\label{eq:fixed-flux-ratio}
 \frac{\mu(x)P_m(x,y)}{\mu(y)P_{\bar m}(y,x)}
 =\frac{\theta_m}{\theta_{\bar m}}.
\end{equation}
The prepared marked rows recover \(\theta_m\) by summation and recover
\(C_m\) by division and multiplication by \(\mu(x)\).

A profile records the positive marked support, a nested schedule of
partitions \(F_0,\ldots,F_T\) of \(X\) with fixed initial partition
\(F_0\), a route license, and a set of citeable checks from a fixed
finite catalogue. At step \(t\) the probability of each mark and next
block of \(F_{t+1}\) must be constant on each current block of \(F_t\),
a stepwise lumpability condition; compare the block-sum
criterion~\citep[Theorem 6.3.2]{kemenysnell}.
Decisions at grade \(t\) are binary functions on the blocks of \(F_t\).
The route license is a congruence on the monoid of mark words up to
a fixed length, with an absorbing overflow element. An identified
pair must have equal length, equal probability, and equal conditional
projected endpoint laws. A live check specifies a positive rational
ratio of probabilities of two words of equal length. Each law
holding the check supplies an exact executed replay bound to that
law. Historical checks refer to immutable source records.
Profiles are attained if some lawful rational law
realizes them. The order on attained profiles contracts support,
refines partitions, and enlarges the congruence and check set.

Let \(H_F\) be the subspace of \(L^2(\mu)\) of functions constant on
the blocks of \(F\). The stepwise disclosure tests are equivalent to
\begin{equation}\label{eq:fixed-terminal}
 K_m H_{F_T}\subseteq H_{F_0}\qquad(m\in M).
\end{equation}
Indeed, the stepwise conditions are \(K_m H_{F_{t+1}}\subseteq H_{F_t}\).
By nesting, \(H_{F_t}\) is invariant under every \(K_m\) for \(t<T\).
Since \(K_m^*=K_{\bar m}\) by \eqref{eq:fixed-flux}, its orthogonal
complement is also invariant. Hence \(K_m\) annihilates
\(H_{F_{t+1}}\cap H_{F_t}^{\perp}\). Summing these orthogonal increments
gives \eqref{eq:fixed-terminal}, which in turn implies each stepwise
condition by nesting. The observed marked paths preserve the full
stationary path-reversal likelihood ratio. The micro path ratio is
the product of \eqref{eq:fixed-flux-ratio} over its marks, hence is
constant on each observed path class. Summing over that class gives
the same ratio against the reversed path with reversed observation
schedule and marks.

For a word \(w\) with count vector \(a(w)\), its probability from any
preparation is \(\theta^{a(w)}\) and its conditional projected kernel
is the product of the \(Q_m\) in word order. Equal summaries stay
equal in common contexts by multiplication of probabilities and
kernels, and equal lengths overflow together. Thus on an attained
profile the nontrivial rate constraints from the licenses and live
checks have the form
\begin{equation}\label{eq:fixed-rate-rows}
 \theta^{d_i}=c_i,\qquad
 d_i\in\mathbb Z^M,\quad \sum_{m\in M}(d_i)_m=0,\quad c_i\in\QQ_{>0}.
\end{equation}
Write \(D\) for their exponent matrix. The conditional-flux tests
are rational affine equations by \eqref{eq:fixed-flux} and
\eqref{eq:fixed-terminal}, with positivity on the profile support
and zero elsewhere. The law fiber therefore separates into
conditional-flux and rate factors.

An elementary capability edit is a cover of attained profiles.
A support cover keeps the rates fixed, and the other covers keep
the full law fixed. Numerical edits hold the profile fixed. For
fluxes, irreducible numerical directions are circuits of the
homogeneous equations on the fixed support, meaning nonzero
directions of inclusion-minimal support, with collinear
subdivisions counted as one direction. These are the flux changes
that admit no noncollinear conformal splitting, the elementary vectors
of Rockafellar~\citep{rockafellar}. For rates, use
the additive coordinates
\(v_p(\theta_m/\theta_n)\) for primes \(p\) and ordered pairs
\((m,n)\), where \(v_p\) is the prime exponent of a positive
rational. An irreducible rate change is supported at one prime and
its difference cannot be split into two nonzero lawful differences
conformal to it. A difference is conformal to another if each
nonzero entry has the same sign and no greater magnitude. The
irreducible rate changes at a prime are thus the Graver basis
elements of the lattice of lawful prime-exponent displacements~\citep{graver}.

\begin{theorem}[Fixed-projection normal form]\label{thm:fixed-normal-form}
Two lawful rational endpoints are connected by these edits exactly
when their profiles are ordered in the direction of the edit.
The capability covers each change a single capability. Disclosure
covers split one block at one grade, route covers adjoin one pair
and take congruence closure, and publication covers add one check.
Together with the numerical changes, the effects are the six
in \Cref{tab:effects}, distinguished by the same fields and clock
grades. Test-preserving relabelings preserve the effect types and
the elementary numerical directions.
\end{theorem}

\begin{proof}
For a feasible support contraction from \(C\) to a retained mask
\(N\), optimize over the closed polytope of conditional-flux tables
\(Z\) satisfying the old affine tests and vanishing off \(N\).
Maximize \(\eta\geq0\) with \(Z_e\geq\eta C_e\) for every \(e\in N\).
A feasible table positive on \(N\) implies a positive optimum.
The objective is bounded by \(|M|/\sum_{e\in N}C_e\), and a rational
vertex attains it. On its optimal face minimize
\(\sum_{e\in N}(Z_e-C_e)^2/C_e\). On the smallest face containing
the minimizer, the affine equations of the face and the linear
gradient conditions along its tangent space have rational
coefficients. They determine the minimizer uniquely by strict
convexity, hence it is rational. Its support is \(N\). Keeping the
rates fixed preserves every old rate test.

At a fixed profile let \(u\) be the target conditional-flux
displacement. Choose a rational homogeneous direction \(v\ne0\)
with signs compatible with \(u\), supported in \(\supp u\), and
with inclusion-minimal support among such choices. If a nonzero
homogeneous direction were supported on a proper subset of
\(\supp v\), one could add a rational multiple in one of its two
directions up to the first zero of \(v\), keeping the signs and
leaving a nonzero result. This contradicts minimality. Thus \(v\)
is a circuit. The direction space on its support is one-dimensional
because two independent directions on that support would have a
nonzero linear combination canceling a coordinate. Thus any
conformal splitting of a circuit is collinear. Scale \(v\) by
\(\min_{e\in\supp v}|u_e/v_e|\) and subtract that step from the
residual \(u\). At least one more entry is at its target and every
entry stays between its original and final values. The affine
tests and positivity on the support are preserved. For a
displacement of nonminimal support the first step is supported on
a proper subset of \(\supp u\), giving a noncollinear splitting.
Iteration terminates because the residual support strictly decreases.

For the rate factor, a pair of lawful rate vectors
\(\theta,\widehat\theta\) has
\(z_p(m)=v_p(\widehat\theta_m/\theta_m)\) in \(\ker D\) at each prime,
by unique factorization in \eqref{eq:fixed-rate-rows}. Only finitely
many primes occur. The possible displacements at a prime are the
integer lattice
\[
 \{(z_m-z_n)_{m,n\in M}:z\in\mathbb Z^M,\ Dz=0\}.
\]
Choose a nonzero lattice element of least \(\ell^1\)-norm among
those conformal to the residual at that prime. It is irreducible,
since any conformal splitting would give a smaller such element.
Subtract and repeat until the integer norm is zero. The element
with lift \(z\) is realized by
\begin{equation}\label{eq:fixed-prime-move}
 \theta'_m=\frac{\theta_m p^{z_m}}{\sum_{n\in M}\theta_n p^{z_n}}.
\end{equation}
Normalization cancels in the homogeneous tests and pairwise rate
ratios. This jump changes only the prescribed prime
in the ratio coordinates. Matching all pairwise ratios gives the
target normalized rates.

If the endpoint profiles are ordered, the target law also attains
the profile with the target mask and the source's other fields.
This follows from block summation in \eqref{eq:fixed-terminal} and
inclusion of the held tests. Every support-changing cover therefore
changes only the support. If the support agrees, the upper law
attains each intermediate choice of nested partitions, congruence
and checks between the two profiles. Thus a cover can change only
one of those fields. Among different grades, advancing the latest
differing one preserves nesting. Its blocks can be split one at a
time into two unions of upper blocks, so a cover makes just one
binary split. For a congruence cover, adjoining any new pair and
taking its congruence closure stays below the upper congruence.
The cover forces equality. A check cover adds one name.

Take successive support covers first, preserving the earlier tests.
Apply the circuit and prime decompositions on the final mask, then
take the remaining capability covers at the final law. The attained
order is finite and the numerical decompositions terminate. Every
allowed step is monotone in the profile order, proving the converse.
The readouts recover the masks, partitions, congruences, inventories,
conditional fluxes and rate ratios used in these tests. Relabeling
preserves the homogeneous spaces and support of flux differences,
and permutes the pairwise rate coordinates, so it preserves the
elementary numerical directions and capability effects. The example
below supplies the joint lower bound.
\end{proof}

\paragraph{A coupled edit with a changing affinity.}
Take \(B=\mathbb Z/4\) and \(X=B\times\{0,1,2,3\}\) with uniform
\(\mu\). The marks \(+,-,c,h\) translate the base by \(1,-1,2,0\),
respectively, with \(+\) and \(-\) reversed and the other marks
self-reversing. Initially every conditional register table resets
uniformly and the partitions at grades \(0,1,2\) are indiscrete.
Use the length-two route horizon. Normalize the initial rate weights
\((4,1,3,2)\), hold \(+-\sim hh\), and hold
\(\theta_c/\theta_-=3\). Both licensed words have zero displacement
and equal probability, and all nonempty contexts for them overflow.

At base point \(0\) for mark \(h\), delete the two register
transitions between \(0\) and \(1\) and add the deleted entries to
the two corresponding diagonals. This is a support cover because
the deleted proposals form one reversal pair. For the rational
gate in the proof, untouched rows bound the optimum by \(\eta=1\).
The stated table attains this bound. At equality the row and
column sums leave only the indicated diagonal increases. In the
clockwise table from base point \(0\), change the conditional register fluxes at
\((0,2),(1,3),(0,3),(1,2)\) by
\((1,1,-1,-1)/32\), with transposed changes in its reverse table.
Here conditional register flux uses the uniform register weight
\(1/4\). The zero-margin space on these entries and their reverses
is one-dimensional and any nonzero direction uses all of them.
This is a circuit and makes the register law depend on base position.

Apply \eqref{eq:fixed-prime-move} with lifts \((-2,0,0,-1)\) at
\(2\) and \((2,0,0,1)\) at \(3\). The successive rate weights are
proportional to \((1,1,3,1)\) and \((9,1,3,3)\). The equations
\(\theta_+\theta_-=\theta_h^2\) and \(\theta_c/\theta_-=3\) hold
throughout. The relative lattice has rank one and each move
changes \(v_p(\theta_h/\theta_-)\) by one in magnitude, so both are
irreducible. The clockwise four-cycle affinity changes from \(4^4\)
to \(9^4\).

\subsection{Which affinities the held tests determine}

\begin{theorem}[Affinity row-space criterion]\label{thm:fixed-affinity}
At a fixed attained profile, a zero-sum integer exponent
\(d\in\mathbb Z^M\) gives an invariant relative monomial
\(\theta^d\) under lawful numerical edits exactly when \(d\) is in
the rational row space of \(D\). Outside that space an exact rational
edit chain starting at any lawful rate vector changes \(\theta^d\).
The real rate fiber in the open positive simplex is smooth of
dimension \(|M|-1-\operatorname{rank}D\). For a fixed supported
closed marked walk, its multiplicative affinity is the monomial
with \(d=\sum_j(e_{m_j}-e_{\bar m_j})\).
\end{theorem}

\begin{proof}
The centered log coordinates
\(u_m=\log\theta_m-|M|^{-1}\sum_n\log\theta_n\)
identify the open positive simplex with the zero-sum hyperplane.
The inverse exponentiates and normalizes. The held equations are
\(Du=\log c\), with all rows in the zero-sum hyperplane, giving the
dimension. A row combination makes \(d\cdot u\) constant. Otherwise
rational elimination and clearing denominators give an integer
\(z\) with \(Dz=0\) and \(d\cdot z\ne0\). For any prime \(p\),
\eqref{eq:fixed-prime-move} preserves the held equations and changes
the monomial by \(p^{d\cdot z}\). The normal form realizes this
change by lawful rational numerical steps with the conditional
tables fixed. Multiplying \eqref{eq:fixed-flux-ratio} along the walk
gives the affinity formula.
\end{proof}

In the real fixed-profile fiber the rate factors are connected by
centered log segments and the conditional-flux factors by affine
segments on the fixed support. The proof below changes the permitted
probability tests and exhibits fixed profiles with isolated laws.

\subsection{Finite-event comparisons and failure of locality}

\begin{theorem}[Finite-event expressivity]\label{thm:finite-event}
For every rational polynomial on a positive rate simplex there are
two fixed events of the same finite duration in a one-state marked
Markov instrument whose probabilities agree exactly on the
polynomial's zero set. In particular, for any \(k\geq2\) one such
test gives a fixed-profile real law fiber consisting of exactly
\(k\) isolated rational laws. Below a positive native tolerance,
no numerical edit chain connects distinct laws in this fiber.
\end{theorem}

\begin{proof}
Clear denominators and homogenize to degree \(d\) using
\(\sum_m\theta_m=1\). The polynomial on the simplex is then
\(\sum_a n_a\theta^a\) for integers \(n_a\). Use a one-state
instrument with independent marks of rates \(\theta_m\), all
self-reversing. A constant truth value can be realized with two
identical events or with an empty and a full event, including on
a one-point simplex.

For at least two marks, choose distinct \(i,j\) and an integer
\(r\geq1\) with \(\binom{2r}{r}\geq\max_a |n_a|\).
There are that many words containing \(r\) copies each of \(i,j\),
all of probability \((\theta_i\theta_j)^r\). For each \(a\) with
\(n_a\ne0\), choose one length-\(d\) prefix with multiplicities \(a\)
and append \(|n_a|\) distinct such suffixes. The resulting words
are all distinct and have length \(d+2r\). Put the words for
positive coefficients in one event and those for negative
coefficients in the other. Their probability difference is
\[
 (\theta_i\theta_j)^r\sum_a n_a\theta^a.
\]
The prefactor is strictly positive. Common prefix and suffix
contexts multiply both event probabilities by the same factor
by independence, so the comparison also supports substitution
in those contexts.

With two rates \(t,1-t\), take the polynomial
\(\prod_{i=1}^k(t-a_i)\) for distinct rational
\(a_i\in(0,1)\). The support, disclosure schedule and held test
are identical at the \(k\) resulting laws. They are all the lawful
rates at this profile. The least distance \(\min_{i\ne j}|a_i-a_j|\)
in the native probability metric is positive, so smaller steps
cannot move between them. The same construction works under a
common rate floor \(0<\delta<1/2\) by choosing the roots in
\((\delta,1-\delta)\).
\end{proof}
